\documentclass[journal=ancac3,manuscript=article,layout=twocolumn]{achemso}
\usepackage{graphicx}\usepackage{amsmath}\usepackage{amssymb}\usepackage{booktabs}\usepackage{caption}
\graphicspath{%
 {D:/DINOSR/Claude/PaperRun_claude/dino_sr_contrastive/docs/explainer/figs/}%
 {D:/DINOSR/Claude/PaperRun_claude/dino_sr_contrastive/docs/interpretation_reports/IMC_SI5__m0.9700_seed42_K60/}%
 {D:/DINOSR/Claude/PaperRun_claude/dino_sr_contrastive/docs/interpretation_reports/EuInAs_B100__m0.9700_seed42_K60/}%
}
\author{Daniel Khaykelson}
\email{author@institution.edu}
\affiliation{Department of Materials Science and Engineering, [Institution]}
\title{Label-Free, Bias-Free Classification of Difficult 4D-STEM Diffraction
with a Diffraction-Tailored Self-Supervised Model}
\keywords{4D-STEM, nanobeam electron diffraction, self-supervised learning,
unsupervised clustering, low-dose, interpretability, indomethacin}
\begin{document}

\begin{abstract}
Four-dimensional scanning transmission electron microscopy (4D-STEM) of
beam-sensitive, low-dose, poorly-ordered materials produces diffraction data
that are hard to segment into meaningful regions without labels, a forward
model, or subjective choices such as the number of clusters or components. We
introduce a self-supervised classifier purpose-built for nanobeam electron
diffraction (NBED): a truncated convolutional DINO encoder trained with
azimuthal ($\theta$-roll) augmentation for in-plane-rotation invariance, a
radial-gated supervised-contrastive objective, and a one-dimensional radial
clustering loss, with emergent cluster count and optional pairwise
semi-supervision. The model needs no labels, no preset $K$, and no
per-feature thresholds. On a sodium-phyllosilicate (NaPHI) flake it reproduces
our previously published, expert-derived domain segmentation label-free and
without the slow SAM + orientation-mapping pipeline (Dice $0.81$, IoU $0.68$;
$\sim$75\% of the indexed domain recovered), and transfers to a new
acquisition in milliseconds. A built-in interpretability protocol
(representational probing, causal ablation, classical-baseline comparison, and
ACOM cross-check) shows the clusters are governed by scattered intensity and
low-$q$ two-dimensional diffraction structure, are rotation-invariant, and
quantifies when the model is distinctive versus classical analysis: highly
distinctive on hard low-dose molecular films (indomethacin, IMC) where
classical decompositions cannot produce a coherent map, and convergent with
classical analysis on well-ordered samples (NaPHI, EuInAs). The method is most
valuable exactly where existing tools fail.
\end{abstract}

\section{Introduction}
4D-STEM records a full diffraction pattern at every probe position, encoding
local structure, orientation, strain and disorder across a
field of view.\cite{ophus2019} Turning that data cube into a meaningful map of
distinct regions is a prerequisite for almost any downstream analysis. For
robust, well-ordered crystals this is routine; for the beam-sensitive,
low-dose, partially-ordered materials of greatest current interest---organic
semiconductors, pharmaceutical polymorphs, layered and defective
crystals---it is not. Patterns are sparse and noisy, domains are numerous,
and diffuse scattering blurs the distinctions a human or a template would key
on.

Unsupervised and self-supervised analysis is attractive because it needs no
labels.\cite{caron2021,uhlemann2023} But the standard automated routes each
demand a subjective choice the microscopist must defend: $k$-means and density
clustering require the number of clusters (or a density scale); non-negative
matrix factorization (NMF) and principal component analysis (PCA) require the
number of components;\cite{spurgeon2021} supervised segmentation requires
labels, which on hard data is prohibitive; and segmentation-model pipelines
(e.g.\ SAM) plus orientation mapping are slow and parameterized. Just as
importantly, even when a clustering succeeds, its \emph{physical} meaning is
rarely established---so the map cannot be trusted.

Here we address both problems. We present a self-supervised classifier
engineered specifically for NBED, which produces a label-free, threshold-free
first segmentation; we validate it against an independent, previously
published expert result; and we pair it with an interpretability protocol that
establishes what the clusters physically represent and when the method adds
value over classical analysis.

\section{Results and Discussion}

\subsection{A diffraction-tailored self-supervised model}
Our classifier is a self-distillation (DINO) network,\cite{caron2021} but every
major design choice is dictated by the physics of NBED rather than inherited
from natural-image vision:
\begin{itemize}
\item \textbf{Convolutional, truncated encoder.} We use a truncated ResNet
rather than a vision transformer. NBED patterns are small, sparse and
low-count; the convolutional inductive bias and the reduced depth match the
available signal, whereas transformers are data-hungry and over-parameterized
for this regime.
\item \textbf{Physics-chosen augmentations.} An azimuthal ``$\theta$-roll''
augmentation rotates the pattern in-plane during training, enforcing
invariance to in-plane crystal rotation---a diffraction pattern rotated about
the optic axis represents the \emph{same} structure. Flips and colour-jitter,
standard in vision, are disabled because they are meaningless or harmful for
diffraction.
\item \textbf{Radial-gated contrastive + 1-D loss.} A radial-gated
supervised-contrastive term\cite{khosla2020} and a one-dimensional radial
clustering loss bias grouping toward the physically meaningful radial/ring
signature.
\item \textbf{Emergent $K$ and optional semi-supervision.} We over-specify the
number of prototypes; the model populates only those it needs, so the active
cluster count emerges from the data rather than being preset. A small number
of pairwise ``same/different'' hints can optionally steer the partition
without reverting to full labelling.
\end{itemize}
The result is a single forward pass from raw cube to a class map, with no
user-set $K$, components, or thresholds.

\subsection{Label-free concordance with expert segmentation}
We first establish trust on a sample with an independent, published answer. In
prior work on a NaPHI flake we segmented an oriented ``Line'' domain using a
SAM segmentation plus orientation mapping with user-defined line
parameters.\cite{ourNaPHI} Running the present model on the same scan, with no
labels and no line parameters, the union of two emergent classes reproduces
that domain (Figure~\ref{fig:naphi}): Dice $=0.81$, IoU $=0.68$, with
$\sim$75\% of the published domain recovered and $\sim$89\% of the model's
domain falling within the published one. The model achieves this without the
SAM step, and a model trained on this flake transfers to a second, tilted
acquisition in milliseconds---turning a slow, parameterized pipeline into an
automatic one.

\begin{figure}[t]\centering
\includegraphics[width=\linewidth]{naphi_concordance.png}
\caption*{\small\textbf{Figure 1.} NaPHI: the label-free DINO class map (centre)
versus our previously published SAM+orientation segmentation (left), same scan.
Overlay (right): green = agreement. The model recovers the oriented domain
(Dice 0.81, IoU 0.68) with no labels, no SAM and no line parameters.}
\label{fig:naphi}
\end{figure}

\subsection{Resolving low-dose molecular films where classical methods fail}
The payoff is on the hard target. On low-dose indomethacin (IMC) films the
model produces coherent, fine-grained class maps, whereas classical automated
methods do not. Clustering virtual dark-field, azimuthal-profile, PCA and NMF
features of the \emph{same} data into the same number of groups agrees only
weakly with the DINO map (best adjusted Rand index, ARI $\approx0.11$--$0.21$
across three IMC films; Figure~\ref{fig:imc}). There is no external ground
truth here---classical analysis simply cannot deliver a usable map---so the
method is the only practical route, and its trustworthiness rests on the
validation above and the interpretability below.

\begin{figure*}[t]\centering
\includegraphics[width=0.92\textwidth]{imc_all_classical.png}
\caption*{\small\textbf{Figure 2.} All three IMC films (SI3, SI4, SI5; rows):
the DINO map (far left of each row) versus classical clusterings of the same
data (virtual-DF, azimuthal-profile, PCA, NMF, combination). Colours are
arbitrary; spatial shape is what matters. Across all three, classical methods
recover only coarse regions and do not reproduce the DINO map (best ARI
$0.11$, $0.15$, $0.21$ for SI3/SI4/SI5).}
\label{fig:imc}
\end{figure*}

\subsection{What the clusters encode}
To establish meaning we apply four complementary tests
(Figure~\ref{fig:interp}). \emph{Representational probing}---5-fold
ridge regression from the frozen embedding to candidate physical factors, with
the correlation ratio $\eta^2$---shows the embedding most strongly encodes the
total scattered intensity and crystallinity, and essentially ignores
crystallographic orientation and Bragg-peak count. \emph{Causal ablations}
(re-inference on perturbed inputs, scored by ARI against the original map)
disentangle the contributions. A consistent post-mask intensity normalization
(rescaling the kept, beam-masked region to constant) collapses the
intensity-dominated samples (NaPHI, EuInAs $\to$ 1--2 classes) but leaves IMC
with structure (7 classes, ARI $0.12$): IMC depends on scattered intensity
\emph{and} low-$q$ 2-D structure, not intensity alone. A complementary
\emph{per-ring} normalization---flattening the radial profile while preserving
within-ring angular structure---collapses the molecular samples (NaPHI, IMC;
i.e.\ radial/intensity-driven and rotation-invariant) but EuInAs retains
ARI $0.29$, showing it additionally uses angular/orientation structure,
consistent with its zone-axis correlation (AMI $\approx0.30$). Azimuthal
averaging collapses all samples (the full 2-D pattern is needed, not a 1-D
profile), and masking the low-$q$ band is far more damaging than high-$q$. \emph{Classical-baseline comparison} quantifies novelty (above).
\emph{ACOM cross-check} confirms the partition is uncorrelated with indexed
zone axis on the molecular samples (adjusted mutual information $\approx0$),
consistent with the $\theta$-roll-imposed rotation invariance. Together: the
clusters are driven by scattered intensity and low-$q$ 2-D diffraction
structure, are rotation-invariant, and their distinctiveness from classical
analysis is material-dependent.

\begin{figure*}[t]\centering
\includegraphics[width=0.95\textwidth]{cross_all.png}
\caption*{\small\textbf{Figure 3.} Cross-sample interpretability (IMC SI3/SI4/SI5,
NaPHI, EuInAs). (a) classical agreement with the DINO map (rises from
distinctive IMC to substantially-captured NaPHI/EuInAs); (b) intensity
dependence (low bar = strongly depends); (c) low-$q$ vs high-$q$ (low bar =
that band matters---NaPHI/EuInAs depend almost entirely on low-$q$).}
\label{fig:interp}
\end{figure*}

\subsection{Generality: crystalline, multi-phase, oriented data}
The approach is not limited to molecular films. On a layered EuInAs sample with
multiple phases, the model produces a sensible multi-domain map; classical
methods here recover about half of it (ARI $\approx0.44$), i.e.\ the model is
\emph{least} distinctive on well-ordered material---it is most distinctive
precisely on the hard molecular case. A three-phase ACOM analysis
(EuInAs/In/As) further shows that, for this crystalline sample, the DINO
classes additionally track crystal orientation (zone-axis AMI $\approx0.30$), a
genuine secondary axis that the molecular films lack.

\subsection{A graphical interface for preparation, analysis and comparison}
To make the method usable by experimentalists who do not program, we provide a
graphical application built around three roles.

\textbf{(1) Data preparation.} It loads the common cube formats (.prz/.npy/
.npz/.h5, including detector-linked masters), and offers the preprocessing that
matters for hard data: intensity scaling (linear $v_{\max}$ or logarithmic
stretch), centre masking, real- or reciprocal-space Gaussian binning for
low-dose denoising, region-of-interest selection, and multi-dataset handling
---all non-destructive, writing new derived datasets.

\textbf{(2) Output analysis.} This is the core value. The class map is fully
interactive: clicking a position shows its raw pattern; clicking a class shows
the \emph{class-average} diffraction; grains can be averaged individually; and
single-frame, per-grain and per-cluster averages can be stacked and compared
side by side, with on-image $q$/$d$-spacing readout and gradient-attribution
overlays. A one-click ``Interpretation'' tab runs the full battery above
(probing, ablation, classical-baseline, ACOM cross-check) and writes a report
---so the question ``what do the clusters mean?'' is answered inside the same
tool that produced them.

\textbf{(3) Classical analysis, for comparison or standalone.} The same
interface runs the classical methods on the same map---automated crystal
orientation mapping (ACOM), NMF + clustering, crystallinity and
spottiness mapping, and SAM-based segmentation---so a user can directly compare
the self-supervised map against any of them, or use the classical tools on
their own. The model is thus offered not as a black-box replacement but as the
fast, unbiased first pass within a complete, comparison-ready toolkit.

\section{Conclusions}
We have introduced a self-supervised, diffraction-tailored classifier that
turns difficult low-dose 4D-STEM into a label-free, threshold-free, and
physically interpretable domain map. It reproduces an independent published
expert segmentation without labels or SAM, transfers across acquisitions in
milliseconds, and---uniquely among the methods we compared---extracts a
coherent map from low-dose molecular films that classical decompositions
cannot. A built-in interpretability protocol establishes that the clusters
encode scattered intensity and low-$q$ diffraction structure rather than
orientation, and quantifies where the method adds value. Because it removes the
subjective choices (labels, $K$, components, thresholds) that gate existing
workflows while remaining auditable, it is well suited as a first-pass,
high-throughput tool for beam-sensitive crystalline materials.

\section{Methods}
\textbf{Model and training.} Truncated ResNet encoder; DINO self-distillation
with teacher/student temperatures and EMA; azimuthal $\theta$-roll augmentation
(student/teacher ranges), flips and colour-jitter disabled; radial-gated
supervised-contrastive term and 1-D radial clustering loss; over-specified
prototype count with emergent active $K$; optional pairwise semi-supervision.
\textbf{Data.} NBED cubes of IMC (low-dose molecular), NaPHI and EuInAs.
\textbf{Interpretability.} Ridge-probe $R^2$ and $\eta^2$; causal input
ablations re-inferred and scored by ARI; classical baselines (virtual-DF,
azimuthal-PCA, pattern-PCA, NMF) clustered to the active $K$ and scored by
ARI/AMI; ACOM via py4DSTEM\cite{savitzky2021,ophus2022}. \textbf{Concordance.}
IoU/Dice/containment of aligned, same-grid class and reference maps.

\begin{acknowledgement}
Analyses were performed on frozen, previously validated model checkpoints.
\end{acknowledgement}

\begin{thebibliography}{9}
\bibitem{ophus2019} Ophus, C. Four-Dimensional Scanning Transmission Electron
Microscopy (4D-STEM). \emph{Microsc.\ Microanal.} \textbf{2019}, \emph{25},
563--582.
\bibitem{caron2021} Caron, M.\ et al. Emerging Properties in Self-Supervised
Vision Transformers. \emph{ICCV} \textbf{2021}, 9650--9660.
\bibitem{uhlemann2023} Self-supervised representation learning for 4D-STEM.
\emph{(representative prior work)}.
\bibitem{spurgeon2021} Spurgeon, S.\ R.\ et al. Towards data-driven
next-generation transmission electron microscopy. \emph{Nat.\ Mater.}
\textbf{2021}, \emph{20}, 274--279.
\bibitem{khosla2020} Khosla, P.\ et al. Supervised Contrastive Learning.
\emph{NeurIPS} \textbf{2020}, \emph{33}, 18661--18673.
\bibitem{ourNaPHI} [Author]\ et al. \emph{Nano Lett.} \textbf{2025}
(NaPHI 4D-STEM study; DOI 10.1021/acs.nanolett.5c04946).
\bibitem{savitzky2021} Savitzky, B.\ H.\ et al. py4DSTEM. \emph{Microsc.\
Microanal.} \textbf{2021}, \emph{27}, 712--743.
\bibitem{ophus2022} Ophus, C.\ et al. Automated Crystal Orientation Mapping.
\emph{Microsc.\ Microanal.} \textbf{2022}, \emph{28}, 390--403.
\end{thebibliography}
\end{document}
